\subsubsection{Kosinussatz}
\label{subsubsec:fundamentals_oblique_triangles_cosine_rule}

Der Kosinussatz verallgemeinert den Satz des Pythagoras auf beliebige
Dreiecke. Neben drei Seitenlängen tritt dabei der eingeschlossene Winkel auf.

\begin{TheoremBox}{Kosinussatz}{fundamentals_cosine_rule}
  In einem beliebigen Dreieck gilt
  \begin{ClozeArea}[27]{5}
    \[
      c^2=a^2+b^2-2ab\cos(\gamma).
    \]
  \end{ClozeArea}
  Dabei liegt der Winkel $\gamma$ der Seite $c$ gegenüber.
\end{TheoremBox}

Da keine Seite ausgezeichnet ist, kann dieselbe Beziehung für jede Seite
formuliert werden:
\begin{ClozeArea}[30]{8}
  \begin{align*}
    a^2&=b^2+c^2-2bc\cos(\alpha),\\
    b^2&=a^2+c^2-2ac\cos(\beta),\\
    c^2&=a^2+b^2-2ab\cos(\gamma).
  \end{align*}
\end{ClozeArea}

\begin{RemarkBox}
  Für $\gamma=90^\circ$ ist $\cos(\gamma)=0$. Damit geht der Kosinussatz in
  den Satz des Pythagoras über:
  \[
    c^2=a^2+b^2-2ab\underbrace{\cos(90^\circ)}_{=0}=a^2+b^2.
  \]
\end{RemarkBox}
